%%%%%%%%%%%%%%%%%%%%%%%%%%%%%%%%%%%%%%%%%
% University Assignment Title Page 
% LaTeX Template
% Version 1.0 (27/12/12)
%
% This template has been downloaded from:
% http://www.LaTeXTemplates.com
%
% Original author:
% WikiBooks (http://en.wikibooks.org/wiki/LaTeX/Title_Creation)
%
% License:
% CC BY-NC-SA 3.0 (http://creativecommons.org/licenses/by-nc-sa/3.0/)
% 
% Instructions for using this template:
% This title page is capable of being compiled as is. This is not useful for 
% including it in another document. To do this, you have two options: 
%
% 1) Copy/paste everything between \begin{document} and \end{document} 
% starting at \begin{titlepage} and paste this into another LaTeX file where you 
% want your title page.
% OR
% 2) Remove everything outside the \begin{titlepage} and \end{titlepage} and 
% move this file to the same directory as the LaTeX file you wish to add it to. 
% Then add \input{./title_page_1.tex} to your LaTeX file where you want your
% title page.
%
%%%%%%%%%%%%%%%%%%%%%%%%%%%%%%%%%%%%%%%%%
%\title{Title page with logo}
%----------------------------------------------------------------------------------------
%	PACKAGES AND OTHER DOCUMENT CONFIGURATIONS
%----------------------------------------------------------------------------------------

\documentclass[12pt]{article}

\usepackage[francais]{babel}
\usepackage[utf8x]{inputenc}
\usepackage[T1]{fontenc}
\usepackage{color}

\usepackage{amsmath}
\usepackage{graphicx}
\usepackage{enumerate}
\usepackage{url}

% Define new command
\newcommand{\HRule}{\rule{\linewidth}{0.5mm}}

\newcommand{\crt}{\emph{Nexys 4 DDR\ }}
%\def\thesubsection{\alph{section}}

\begin{document}

\begin{titlepage}

\center % Center everything on the page
 
%----------------------------------------------------------------------------------------
%	HEADING SECTIONS
%----------------------------------------------------------------------------------------

\textsc{\LARGE Universit\'e Pierre et Marie Curie}\\[1.5cm] % Name of your university/college
\textsc{\Large PSESI}\\[0.5cm] % Major heading such as course name

%----------------------------------------------------------------------------------------
%	TITLE SECTION
%----------------------------------------------------------------------------------------

\HRule \\[0.4cm]
{ \huge \bfseries Projet Centrale DCC}\\[0.4cm] % Title of your document
{ \huge \bfseries Document Définition Interface et Logique Enclenchement}\\[0.4cm] % Title of your document
\HRule \\[1.5cm]
 
%----------------------------------------------------------------------------------------
%	AUTHOR SECTION
%----------------------------------------------------------------------------------------

\begin{minipage}{0.4\textwidth}
\begin{flushleft} \large
\emph{\'Etudiant:}\\
Maxime \textsc{AYRAULT} 3203694 % Your name
\end{flushleft}
\end{minipage}
~
\begin{minipage}{0.4\textwidth}
\begin{flushright} \large
\emph{Encadrant:} \\
Julien \textsc{DENOULET} % Supervisor's Name
\end{flushright}
\end{minipage}\\[2cm]

%----------------------------------------------------------------------------------------
%	DATE SECTION
%----------------------------------------------------------------------------------------

{\large \today}\\[2cm] % Date, change the \today to a set date if you want to be precise

%----------------------------------------------------------------------------------------
%	LOGO SECTION
%----------------------------------------------------------------------------------------

%%\begin{figure}
%%  \subfigure[]{\includegraphics[scale=0.2\textwidth]{logo.png}} 
%%\end{figure} 
\includegraphics[width=0.2\textwidth]{logo.png}

%----------------------------------------------------------------------------------------

\vfill % Fill the rest of the page with whitespace

\end{titlepage}



%\begin{abstract}
%Your abstract here.
%\end{abstract}

\section{Introduction}
\label{sec:introduction}
Présentation des 3 interfaces~:
\begin{itemize}
\item Interface Homme Machine
\item Interface Centrale DCC et le circcuit ferroviaire
\item Interface interne Centrale DCC - Moteur IXL
\end{itemize}

et présentation des règles de la logique d'enclenchement.


\section{Interface Homme Machine}
\label{sec:int-dcc}

\subsection{\underline{Interface commande train}}
\label{sec:ihm-train}

bla bla bla...

\subsection{\underline{Interface itinéraire}}
\label{sec:ihm_iti}

bla bla bla...



\section{Interface Centrale DCC - Train, aiguillage et capteurs}
\label{sec:int_dcc}

\subsection{\underline{Interface Train}}
\label{sec:int_train}

bla bla bla...-+

\subsection{\underline{Interface Aiguillage}}
\label{sec:int_aig}

bla bla bla...

\subsection{\underline{Interface Capteurs}}
\label{sec:int_cap}

bla bla bla...


\newpage
\section{\underline{Interface Centrale DCC - Moteur IXL}}
\label{sec:int_ixl}   

La logique d'enclenchement est basée sur les notions de signaux, de circuit de voie et d'itinéraire ferroviaire.

Un \emph{signal} est le moyen de communication entre le conducteur du train et la logique d'enclenchement. 
Il permet d'indiquer au conducteur l'état de la voie se trouvant devant le train (par définition un 
train ne recule jamais, il est nécessaire de changer de cabine de conduite pour changer
la direction du train). Il existe beaucoup de signaux ferroviaire. Les trois plus importants sont~:
\begin{itemize}
\item signal \emph{rouge} qui signifie que le train doit s'arrêter devant le signal car la
voie devant lui n'est pas sûre (présence d'un autre train, aiguillage en train d'être commandé, panne 
matériel...)
\item signal \emph{jaune} qui signifie qu'il y a un danger à proximité. Le train peut passer le signal
mais à vitesse réduite seulement et doit être capable de s'arrêter sur une courte distance. Ceci permet 
de fluidifier le traffic.
\item signal \emph{vert} qui signifie que la voie devant le train est libre de tout danger et que le train 
peut poursuivre sa route en respectant les vitesses limites.
\end{itemize}

L'objectif des opérateurs consiste à toujours présenter un signal vert devant un train de façon 
à ce que le train n'est jamais à s'arrêter à cause d'un danger.


\medskip
Le projet implémente le signal rouge et le signal vert. Le signal jaune ne peut pas être implémenté sans conducteur 
dans le train. Il n'y a pas de signaux sur le circuit ferroviaire, mais la notion de signaux sera quand même
implémentée. Un signal \emph{virtuel} sera créé pour chaque circuit de voie. Il autorisera ou pas l'entrée d'un train dans 
le circuit de voie auquel il est associé.
Ceci permettra de commander les trains. 

\medskip
Un \emph{Ciruit de Voie} (CdV) représente la plus petite partie de voie sur lequel il
est possible de détecter la présence d'un train. Pour un train réel, la détection de la 
présence d'un train s'effectue par les essieux du train. En effet, les essieux du train sont conducteur
(métal), ils peuvent servir comme un \emph{interupteur} entre les 2 rails de la voie. Si une partie du 
train est présent alors les 2 rails sont au même potentiel (le circuit est dit fermé), sinon les rails sont 
a des potentiels différents (le circuit est dit ouvert). 

\medskip
Pour le projet, un Circuit de voie sera délimité par 2 capteurs (sauf en zone d'aiguille où le circuit de voie sera 
délimité par les capteurs de chaque branche). L'état d'occupation du CdV sera fourni pas le sens de parcours 
identifié par le capteur pour chaque train.

\medskip
Un \emph{itinéraire ferroviaire} représente une portion de voie entre 2 signaux. Il contient les mouvements
d’aiguilles nécessaires pour que le train parcourt le chemin entre le signal de
départ et le signal de sortie de l'itinéraire. Dans le système d'enclenchement, un itinéraire
est représenté par une machine à états. Les états possibles sont:
\begin{itemize}
\item \emph{commandé}. L’opérateur demande au train de parcourir l’itinéraire
choisi. Il a par conséquent demandé la formation de l’itinéraire au
système d’enclenchement. Si les conditions sur les aiguilles à manoeuvrer
pour former l’itinéraire sont sûres, le système d’enclenchement
passe à l’état de formation de l’itinéraire sinon il est rejeté.
\item \emph{formé}. Les aiguilles sont positionnées dans la bonne direction. Il n’est
plus possible de bouger les aiguilles tant que le train se situe dans une
certaine zone.
\item \emph{détruit}. Le train a parcouru l’itinéraire (il a passé le signal de sortie), l'itinéraire peut
 être détruit. Ou l'opérateur a demandé la destruction de l'itinéraire (en cas d'erreur par exemple)
Dans ce cas, si le train ne se trouve pas trop proche du signal d'entrée alors l'itinéraire peut 
être détruit sinon la demande de destruction est rejetée. 
Ceci permet d'autoriser le parcours sur les aiguilles et la zone d’un autre train. 
\end{itemize}


L'interface entre la Centrale DCC et le Moteur IXL est bi-directionnelle. 
Dans le sens Centrale vers IXL, les informations envoyées sont~:
\begin{itemize}
\item les commandes d'itinéraire
\item l'état d'occupation des circuits de voie (libre/occupé)
\item l'état des aiguilles (gauche/droite)
\end{itemize}

Dans le sens IXL vers Centrale, les informations envoyées sont~:
\begin{itemize}
\item les commandes des aiguilles
\item l'état des itinéraires (commandé/formé/détruit)
\item l'état des signaux (rouge/vert)
\end{itemize}

\subsection{\underline{Commande itinéraire}}
\label{sec:ixl_iti}

bla bla bla...

\subsection{\underline{\'Etat des Circuits de voie}}
\label{sec:ixl_cdv}

bla bla bla...

\subsection{\underline{\'Etat des aiguillages}}
\label{sec:ixl_aig}

bla bla bla...

\subsection{\underline{\'Etat des itinéraires}}
\label{sec:st_iti}

bla bla bla...


\subsection{\underline{Commande des aiguillages}}
\label{sec:cmd_aig}

bla bla bla...

\subsection{\underline{\'Etat des signaux}}
\label{sec:st_sig}

bla bla bla...
  







\newpage
\section{\underline{Logique Enclenchement}}
\label{sec:log_enc}

\subsection{\underline{Circuit de Voie}}
\label{sec:CdV}

Un Circuit de Voie est la section de voie se trouvant entre 2 capteurs
consécutifs. Les capteurs de présence ne permettent pas de savoir dans
quel sens va le train. Ils permettent uniquement de lire un code barre
placé sous le train.  

\subsection{\underline{Zone d'aiguille}}
\label{sec:aig}

Bla bla bla...


\subsection{\underline{Itinéraire}}
\label{sec:iti}

Bla bla bla...

\subsection{\underline{zone d'aiguille}}
\label{sec:aig}

Bla bla bla...



\newpage

\bibliographystyle{plain}
\bibliography{biblio}

\end{document}
